% ============================================================
%  radicales_quimicaorganica.tex
%  Guía / Key Sheet — Radicales y Grupos Funcionales
%  Química Orgánica · Mathwizards Consultoría Educativa STEM
%
%  Compilar con XeLaTeX:
%    cd tex_files/radicales_quimicaorganica && latexmk -xelatex radicales_quimicaorganica.tex
% ============================================================
\documentclass[12pt, letterpaper]{article}

% ── Paquetes ─────────────────────────────────────────────────
\usepackage{mathwizards-guia}
\usepackage{chemfig}

% ── Configuración de chemfig ─────────────────────────────────
\setchemfig{
  atom sep        = 1.35em,
  bond offset     = 0.12em,
  double bond sep = 0.22em,
  bond style      = {line width = 0.6pt},
  angle increment = 30
}

% ── Estilo de enlace ondulado (valencia libre) ───────────────
\tikzset{
  wavy/.style={decorate, decoration={snake, amplitude=0.8pt, segment length=3.5pt}}
}

\newcommand{\wavybond}{%
  \tikz[baseline=-0.5ex]{%
    \draw[line width=0.6pt, decorate, decoration={snake, amplitude=0.8pt, segment length=3.5pt}]
      (0,0) -- (0.5,0);%
  }%
}

% ── Datos de la guía ─────────────────────────────────────────
\mwguia{Guía — Radicales y Grupos Funcionales}{Ing. Gabriel Astudillo $\cdot$ Química Orgánica}

\begin{document}
% ============================================================

% ── Portada ──────────────────────────────────────────────────
\mwportada{Radicales y Grupos Funcionales}{Key Sheet de Química Orgánica}{Nomenclatura IUPAC y común · Estructuras esqueléticas}

\begin{cuadroinfo}
\textbf{¿Qué encontrarás aquí?} Esta guía reúne, en dos secciones esenciales, las herramientas de nomenclatura y formulación más empleadas en química orgánica:
\begin{itemize}[leftmargin=1.4em, itemsep=2pt, topsep=3pt]
  \item \textbf{Sección 1 — Jerarquía de grupos funcionales:} el orden oficial de prioridad para seleccionar la cadena principal, el sufijo principal y los prefijos cuando actúan como sustituyentes.
  \item \textbf{Sección 2 — Radicales (sustituyentes):} los grupos alquilo, alquenilo, alquinilo, cicloalquilo, arilo y con heteroátomos más frecuentes, con su nombre IUPAC sistemático y común.
\end{itemize}
\end{cuadroinfo}

\vspace{0.6em}

\begin{cuadroteorianobg}
\textbf{Convenciones de representación esquelética empleadas:}
\begin{itemize}[leftmargin=1.4em, itemsep=2pt, topsep=2pt]
  \item \textbf{Línea ondulada (\wavybond):} indica el enlace libre o punto de unión del radical/grupo a la cadena principal.
  \item \textbf{Estructura esquelética (zig-zag):} los vértices e intersecciones representan átomos de carbono con sus respectivos hidrógenos implícitos.
  \item \textbf{Carbonos de los extremos:} únicamente se explicitan los extremos de cadena (grupos terminales como $-\text{CH}_3$, $=\text{CH}_2$, etc.) y heteroátomos ($-\text{OH}$, $-\text{NH}_2$, $-\text{SH}$).
  \item \textbf{Geometría lineal de alquinos:} los enlaces triples ($\text{C}\equiv\text{C}$) mantienen el ángulo estricto de $180^\circ$ característico de la hibridación $sp$.
\end{itemize}
\end{cuadroteorianobg}

\clearpage

% ============================================================
\section{Jerarquía de grupos funcionales}

La tabla siguiente ordena los grupos funcionales de \textbf{mayor a menor prioridad}. El grupo de mayor prioridad determina el \textbf{sufijo} del nombre; los demás se nombran como \textbf{sustituyentes} (prefijos).

\begin{center}
\footnotesize
\renewcommand{\arraystretch}{1.08}
\begin{tabular}{@{}c l c l l@{}}
\toprule
\textbf{N.º} & \textbf{Grupo funcional} & \textbf{Estructura} & \textbf{Sufijo (cadena principal)} & \textbf{Como sustituyente} \\
\midrule
1  & Ácidos carboxílicos      & \chemfig{R-[:30](=[::60]O)(-[::-60]OH)}                              & Ácido \dots-oico            & carboxi- \\
2  & Anhídridos de ácido      & \chemfig{R-[:30](=[::60]O)(-[::-60]O-[:30](=[::60]O)(-[::-60]R))}    & Anhídrido \dots-oico        & — \\
3  & Ésteres                  & \chemfig{R-[:30](=[::60]O)(-[::-60]O-[:30]R)}                        & \dots-oato de alquilo       & alcoxicarbonil- / aciloxi- \\
4  & Halogenuros de acilo     & \chemfig{R-[:30](=[::60]O)(-[::-60]X)}                               & Halogenuro de \dots-oílo    & halocarbonil- \\
5  & Amidas                   & \chemfig{R-[:30](=[::60]O)(-[::-60]NH_2)}                            & \dots-amida                 & carbamoil- \\
6  & Nitrilos                 & \chemfig{R-[:30]~[0]N}                                               & \dots-nitrilo               & ciano- \\
7  & Aldehídos                & \chemfig{R-[:30](=[::60]O)(-[::-60]H)}                               & \dots-al                    & formil- \\
8  & Cetonas                  & \chemfig{R-[:30](=[::60]O)(-[::-60]R)}                               & \dots-ona                   & oxo- \\
9  & Alcoholes                & \chemfig{R-[:30]OH}                                                  & \dots-ol                    & hidroxi- \\
10 & Tioles (mercaptanos)     & \chemfig{R-[:30]SH}                                                  & \dots-tiol                  & mercapto- / sulfanil- \\
11 & Aminas                   & \chemfig{R-[:30]NH_2}                                                & \dots-amina                 & amino- \\
12 & Alquenos                 & \chemfig{R-[:30]=[:-30]-[:30]R}                                      & \dots-eno                   & alquenil- \\
13 & Alquinos                 & \chemfig{R-[:30]~[0]-[:30]R}                                         & \dots-ino                   & alquinil- \\
14 & Éteres                   & \chemfig{R-[:30]O-[:30]R}                                            & éter                        & alcoxi- \\
15 & Sulfuros (tioéteres)     & \chemfig{R-[:30]S-[:30]R}                                            & sulfuro                     & alquiltio- \\
16 & Halogenuros              & \chemfig{R-[:30]X}                                                   & —                           & halógeno- \\
17 & Nitro                    & \chemfig{R-[:30]\chemabove{N}{\scriptstyle +}(=[::60]O)(-[::-60]\chemabove{O}{\scriptstyle -})} & — & nitro- \\
18 & Alcanos                  & \chemfig{R-[:30]CH_3}                                                & \dots-ano                   & alquil- \\
\bottomrule
\end{tabular}
\end{center}

\begin{cuadroextra}
\textbf{Regla práctica:} para nombrar un compuesto polifuncional, identifica el grupo de \textbf{mayor prioridad} (más arriba en la tabla) y úsalo como sufijo principal. Los grupos de menor prioridad se citan como prefijos en orden alfabético. Los grupos 16--18 se nombran siempre como sustituyentes.
\end{cuadroextra}

\clearpage

% ============================================================
\section{Radicales (sustituyentes)}

Un \textbf{radical} es un fragmento molecular que se une a una cadena principal. Se indica su punto de unión mediante la línea ondulada (\wavybond). Las estructuras se ilustran de forma esquelética limpia, explicitando únicamente los carbonos extremos.

\subsection{Radicales alquilo (alcanos)}

\begin{center}
\footnotesize
\renewcommand{\arraystretch}{1.18}
\begin{tabular}{@{}l l c@{}}
\toprule
\textbf{Nombre IUPAC (sistemático)} & \textbf{Nombre común / aceptado} & \textbf{Estructura} \\
\midrule
Metil            & —                & \chemfig{-[0,0.7,,,wavy]CH_3} \\
Etil             & —                & \chemfig{-[0,0.7,,,wavy]-[:30]CH_3} \\
Propil           & $n$-propil       & \chemfig{-[0,0.7,,,wavy]-[:30]-[:-30]CH_3} \\
Butil            & $n$-butil        & \chemfig{-[0,0.7,,,wavy]-[:30]-[:-30]-[:30]CH_3} \\
Pentil           & $n$-pentil       & \chemfig{-[0,0.7,,,wavy]-[:30]-[:-30]-[:30]-[:-30]CH_3} \\
1-metiletil      & Isopropil        & \chemfig{-[0,0.7,,,wavy](-[:50]CH_3)-[:-50]CH_3} \\
2-metilpropil    & Isobutil         & \chemfig{-[0,0.7,,,wavy]-[:30](-[:90]CH_3)-[:-30]CH_3} \\
1-metilpropil    & \emph{sec}-butil ($s$-butil) & \chemfig{H_3C-[:30](-[::90,0.7,,,wavy])-[:-30]-[:30]CH_3} \\
1,1-dimetiletil  & \emph{terc}-butil ($t$-butil) & \chemfig{-[0,0.7,,,wavy](-[:55]CH_3)(-[0]CH_3)-[:-55]CH_3} \\
2,2-dimetilpropil& Neopentil        & \chemfig{-[0,0.7,,,wavy]-[:30](-[:90]CH_3)(-[0]CH_3)-[:-30]CH_3} \\
3-metilbutil     & Isopentil (isoamilo) & \chemfig{-[0,0.7,,,wavy]-[:30]-[:-30](-[:45]CH_3)-[:-75]CH_3} \\
1-metilbutil     & \emph{sec}-pentil ($s$-pentil) & \chemfig{H_3C-[:30](-[::90,0.7,,,wavy])-[:-30]-[:30]-[:-30]CH_3} \\
1,1-dimetilpropil& \emph{terc}-pentil ($t$-pentil)& \chemfig{H_3C-[:30](-[::90,0.7,,,wavy])(-[::-90]CH_3)-[:-30]-[:30]CH_3} \\
\bottomrule
\end{tabular}
\end{center}

\subsection{Radicales alquenilo (alquenos)}

\begin{center}
\footnotesize
\renewcommand{\arraystretch}{1.18}
\begin{tabular}{@{}l l c@{}}
\toprule
\textbf{Nombre IUPAC (sistemático)} & \textbf{Nombre común} & \textbf{Estructura} \\
\midrule
Etenilo          & Vinil        & \chemfig{-[0,0.7,,,wavy]=[:30]CH_2} \\
Prop-2-enilo     & Alilo        & \chemfig{-[0,0.7,,,wavy]-[:30]=[:-30]CH_2} \\
Prop-1-enilo     & Propenil     & \chemfig{-[0,0.7,,,wavy]=[:30]-[:-30]CH_3} \\
2-metilprop-1-enilo & Isobutenil& \chemfig{-[0,0.7,,,wavy]=[:30](-[:80]CH_3)-[:-40]CH_3} \\
But-2-enilo      & Crotil       & \chemfig{-[0,0.7,,,wavy]-[:30]=[:-30]-[:30]CH_3} \\
\bottomrule
\end{tabular}
\end{center}

\clearpage

% ── Fila 1: Alquinilos y Cicloalquilos (side-by-side) ────────
\noindent
\begin{minipage}[t]{0.48\textwidth}
\subsection{Radicales alquinilo}
\begin{center}
\footnotesize
\renewcommand{\arraystretch}{1.15}
\begin{tabular}{@{}l l c@{}}
\toprule
\textbf{IUPAC} & \textbf{Común} & \textbf{Estructura} \\
\midrule
Etinil        & —          & \chemfig{-[0,0.7,,,wavy]~[0]CH} \\
Prop-2-inilo  & Propargilo & \chemfig{-[0,0.7,,,wavy]-[:30]-[:-30]~[0]CH} \\
Prop-1-inilo  & Propinilo  & \chemfig{-[0,0.7,,,wavy]~[0]-[0]CH_3} \\
\bottomrule
\end{tabular}
\end{center}
\end{minipage}\hfill
\begin{minipage}[t]{0.48\textwidth}
\subsection{Radicales cicloalquilo}
\begin{center}
\footnotesize
\renewcommand{\arraystretch}{1.15}
\begin{tabular}{@{}l l c@{}}
\toprule
\textbf{IUPAC} & \textbf{Común} & \textbf{Estructura} \\
\midrule
Ciclopropil   & —          & \chemfig{-[0,0.7,,,wavy]*3(---)} \\
Ciclobutil    & —          & \chemfig{-[0,0.7,,,wavy]*4(----)} \\
Ciclopentil   & —          & \chemfig{-[0,0.7,,,wavy]*5(-----)} \\
Ciclohexil    & —          & \chemfig{-[0,0.7,,,wavy]*6(------)} \\
\bottomrule
\end{tabular}
\end{center}
\end{minipage}

\vspace{0.4em}

% ── Fila 2: Arilos y Heteroátomos (side-by-side) ─────────────
\noindent
\begin{minipage}[t]{0.50\textwidth}
\subsection{Radicales arilo}
\begin{center}
\footnotesize
\renewcommand{\arraystretch}{1.12}
\begin{tabular}{@{}l l c@{}}
\toprule
\textbf{IUPAC} & \textbf{Común} & \textbf{Estructura} \\
\midrule
Fenil        & —               & \chemfig{-[0,0.7,,,wavy]**6(------)} \\
Bencil       & —               & \chemfig{-[0,0.7,,,wavy]-[:30]**6(------)} \\
2-metilfenil & $o$-tolil       & \chemfig{-[0,0.7,,,wavy]**6(-(-[:-120]CH_3)-----)} \\
3-metilfenil & $m$-tolil       & \chemfig{-[0,0.7,,,wavy]**6(--(-[:-60]CH_3)----)} \\
4-metilfenil & $p$-tolil       & \chemfig{-[0,0.7,,,wavy]**6(---(-[0]CH_3)---)} \\
1-naftil     & $\alpha$-naftil & \chemfig{*6((-[::120,0.7,,,wavy])-=*6(-=-=-)-=-=)} \\
2-naftil     & $\beta$-naftil  & \chemfig{-[0,0.7,,,wavy]*6(-=*6(-=-=-)-=-=)} \\
\bottomrule
\end{tabular}
\end{center}
\end{minipage}\hfill
\begin{minipage}[t]{0.46\textwidth}
\subsection{Con heteroátomos}
\begin{center}
\footnotesize
\renewcommand{\arraystretch}{1.12}
\begin{tabular}{@{}l l c@{}}
\toprule
\textbf{IUPAC} & \textbf{Común} & \textbf{Estructura} \\
\midrule
Metoxi   & —        & \chemfig{-[0,0.7,,,wavy]O-[:30]CH_3} \\
Etoxi    & —        & \chemfig{-[0,0.7,,,wavy]O-[:30]-[:-30]CH_3} \\
Propoxi  & —        & \chemfig{-[0,0.7,,,wavy]O-[:30]-[:-30]-[:30]CH_3} \\
Sulfanil & Mercapto & \chemfig{-[0,0.7,,,wavy]SH} \\
Amino    & —        & \chemfig{-[0,0.7,,,wavy]NH_2} \\
Hidroxi  & —        & \chemfig{-[0,0.7,,,wavy]OH} \\
Oxo      & —        & \chemfig{=[0,0.7,,,wavy]O} \\
\bottomrule
\end{tabular}
\end{center}
\end{minipage}

\vspace{0.4em}

\begin{cuadroadvertencia}
\textbf{Ojo con la nomenclatura y el orden alfabético:}
los prefijos \textbf{\emph{sec}-} y \textbf{\emph{terc}-} ($s$-, $t$-) no se consideran para el orden alfabético (se toma la primera letra del grupo base: \emph{sec}-\textbf{b}util entra por la \textbf{b}). En cambio, \textbf{iso-} y \textbf{neo-} forman una sola palabra y \textbf{sí} cuentan para el orden alfabético (\textbf{i}sopropil entra por la \textbf{i}). En nomenclatura IUPAC sistemática estricta se emplean los nombres numerados entre paréntesis (p.~ej. 1,1-dimetiletil).
\end{cuadroadvertencia}

\end{document}